
\subsection*{AI use statement}

(This section is \textbf{required} and does not count toward the page limit.)

In this work, we used generative AI tools for [tasks with required disclosure].
We have not used generative AI tools for [other tasks with required disclosure],
and [the rest of the required disclosure tasks] are not applicable to this work.
Additionally, we used generative AI tools for [tasks with recommended
disclosure]. We have reviewed all AI-assisted work. [Elaborate. For example, “we
checked LLM-generated research ideas for potential plagiarism through a manual
literature survey”, “LLM-generated code was verified and tested for correctness
by 2 authors”, etc.]. We take responsibility for the final content of this work,
including text, claims or artifacts produced with the aid of generative AI.

See the ICLR 2027 AI Policy for Authors for more details. This statement should
not be more than 1 page.

\subsection*{Ethics statement}

(This section is \textbf{recommended} and does not count toward the page limit.)

If authors feel that their paper submission raises questions regarding the Code
of Ethics, they are encouraged to include a paragraph of Ethics Statement (at
the end of the main text before references) to address potential concerns where
appropriate. Topics include, but are not limited to, studies that involve human
subjects, practices to data set releases, potentially harmful insights,
methodologies and applications, potential conflicts of interest and sponsorship,
discrimination/bias/fairness concerns, privacy and security issues, legal
compliance, and research integrity issues (e.g., IRB, documentation, research
ethics). This statement should not be more than 1 page.

\subsection*{Reproducibility statement}

(This section is \textbf{recommended} and does not count toward the page limit.)

It is important that the work published in ICLR is reproducible. Authors are
strongly encouraged to include a paragraph-long Reproducibility Statement at the
end of the main text (before references) to discuss the efforts that have been
made to ensure reproducibility. This paragraph should not itself describe
details needed for reproducing the results, but rather reference the parts of
the main paper, appendix, and supplemental materials that will help with
reproducibility. For example, for novel models or algorithms, a link to an
anonymous downloadable source code can be submitted as supplementary materials;
for theoretical results, clear explanations of any assumptions and a complete
proof of the claims can be included in the appendix; for any datasets used in
the experiments, a complete description of the data processing steps can be
provided in the supplementary materials. Each of the above are examples of
things that can be referenced in the reproducibility statement.


\subsubsection*{Acknowledgments}
Use unnumbered third level headings for the acknowledgments. All
acknowledgments, including those to funding agencies, go at the end of the paper.
